%% kapa: exam=1 round=23 year=2012 period=1 form=A subject=civil_law subject_ko=민법 seq=2 number=2
%% source: Q-Net 감정평가사 기출문제 (cst00309) artlSeq=5052578 file=raw/1008/5052578/2012년 제23회 감정평가사 제1차시험 기출문제(1교시 A형).pdf
%% conversion: kordoc 4.24.51 -> tools/gen1.py (auto); review=proofread-sonnet; flags=none
\begin{question}{2}
기간의 계산에 관한 설명으로 옳지 않은 것은? (다툼이 있으면 판례에 의함)
\begin{choices}
  \item 사원총회를 1주일 전에 통지하여야 하는 경우, 5월 25일 총회 예정일을 기준으로 늦어도 5월 17일 24시까지는 사원들에게 소집통지를 발송하여야 한다.
  \item 정년이 만 60세인 경우, 1952년 7월 8일 오후 3시에 출생한 직원은 2012년 7월 8일 오전 0시에 정년을 맞이한다.
  \item 1981년 5월 20일 오전 11시에 출생한 자는 2001년 5월 21일 오전 0시에 성년으로 된다.
  \item 기간의 말일이 2012년 3월 10일 토요일인 경우에는 2012년 3월 12일 월요일로 만료된다.
  \item 다가오는 2012년 7월 7일부터 1주일까지라고 하면, 2012년 7월 13일 24시에 만료된다.
\end{choices}
\end{question}
\answer{3}
